% ============================================================================
%  Predicting Consequential Volume Errors in Automated Brain-Tumour MRI
%  Segmentation  --  BioAI-2026 (preliminary study, MSD Task01_BrainTumour)
%
%  HOW TO COMPILE (easiest): upload this file to https://www.overleaf.com
%  (New Project -> Upload Project) and click "Recompile". No install needed.
%
%  BEFORE SUBMITTING:
%   1. Fill in the author block (name, supervisor, department, email).
%   2. Add the figure file  risk_coverage.png  next to this .tex
%      (it is your outputs/risk_msd/risk_test_risk_coverage.png).
%   3. Do a final check of references [12]-[14] (2026 preprints).
% ============================================================================
\documentclass[11pt]{article}

\usepackage[utf8]{inputenc}
\usepackage[T1]{fontenc}
\usepackage{times}
\usepackage[margin=1in]{geometry}
\usepackage{graphicx}
\usepackage{booktabs}
\usepackage{amsmath}
\usepackage{amssymb}
\usepackage{siunitx}
\usepackage{hyperref}
\hypersetup{colorlinks=true, linkcolor=blue, citecolor=blue, urlcolor=blue}
\usepackage{caption}
\captionsetup{font=small,labelfont=bf}

\title{\textbf{Predicting Consequential Volume Errors in\\ Automated Brain-Tumour MRI Segmentation}}

\author{
  Subhrakant Sethi\textsuperscript{1} \and Ayush Singh\textsuperscript{1} \and
  Abhayaditya Sogi\textsuperscript{1} \and Pratyaksh Singh\textsuperscript{1}\\[2pt]
  \small \textsuperscript{1}Department of Biotechnology, Thapar Institute of
  Engineering and Technology, Patiala, India\\
  \small Corresponding authors: \texttt{ssethi\_btech24@thapar.edu},
  \texttt{asingh3\_btech24@thapar.edu}
}
\date{}

\begin{document}
\maketitle

% ---------------------------------------------------------------- Abstract
\begin{abstract}
\noindent
Tumour volume measured from brain MRI guides the assessment of glioma progression
and treatment response, yet automated segmentation occasionally mis-estimates
volume enough to mislead clinical interpretation, without signalling that it has
done so. Safe large-scale use therefore requires knowing, for each scan, whether an
output can be trusted. Existing reliability methods largely relate uncertainty to
overlap metrics such as the Dice coefficient rather than to volume error, seldom
define error against a pre-specified clinical threshold, and are rarely framed as an
operational triage task. We study segmentation reliability as the \emph{case-level
prediction of consequential volume error}. Without introducing a new segmentation
architecture or uncertainty method, we take a standard 3D U-Net and ask whether
simple, interpretable models can predict---from the model's own stochastic behaviour
(Monte-Carlo dropout with test-time augmentation)---when an output's whole-tumour
volume error exceeds a clinically motivated threshold ($\geq 25\%$ and $\geq 1$~mL)
fixed before evaluation. Case-level features spanning uncertainty, inter-prediction
disagreement, and boundary/shape are compared, using logistic regression and
gradient boosting, against predicted-volume-only and standard-uncertainty baselines,
with patient-level data splits and a single locked test evaluation. In a preliminary
study on publicly available, expert-annotated MRI, the segmenter reached whole-tumour
Dice $0.81$; reference-free features predicted consequential volume errors with
AUROC~$\approx 0.88$, clearly exceeding the trivial baseline. Their advantage over
standard uncertainty was positive but not statistically significant on the small test
set, except for severe ($\geq 40\%$) errors, where prediction-disagreement
significantly outperformed uncertainty ($p=0.03$). At $90\%$ sensitivity, roughly half
of cases were retainable for automation with about $6\%$ residual error; combining all
features reduced performance. These results provide an honest, reproducible triage
baseline and motivate validation on larger data. This is a research pipeline, not a
clinical decision system.

\vspace{4pt}
\noindent\textbf{Keywords:} brain tumour MRI; segmentation reliability; volumetric
error; uncertainty; quality control; deep learning
\end{abstract}

% ---------------------------------------------------------------- 1 Intro
\section{Introduction}
Gliomas are among the most common and aggressive primary brain tumours, and magnetic
resonance imaging (MRI) is the principal modality for their diagnosis, treatment
planning and longitudinal monitoring~\cite{weller2021}. A central quantitative output
of this imaging is \emph{tumour volume}: changes in the volume of the whole tumour,
tumour core and enhancing region over time are used to assess disease progression and
treatment response, and underpin established response-assessment criteria such as
RANO, which defines progression in part through a $\geq 25\%$ increase in tumour
burden~\cite{wen2010}. Accurate quantification of these volumes is therefore not merely
a segmentation exercise---it directly shapes clinical interpretation.

Deep learning has driven substantial progress in automated brain-tumour segmentation,
with benchmarks such as BraTS establishing 3D convolutional architectures (e.g.,
nnU-Net, SegResNet) as strong, reproducible baselines~\cite{isensee2021,baid2021}. In
principle these models could automate laborious manual contouring. In practice their
outputs cannot be trusted blindly: even high-performing models produce occasional
failures, and a single segmentation that substantially mis-estimates tumour volume can
propagate into an incorrect assessment of progression or response. Deploying automated
segmentation at scale therefore requires more than an accurate model---it requires
knowing, for each individual scan, \emph{whether the resulting output can be trusted},
so that unreliable cases are routed to expert review while the rest are safely
automated.

Prior work has approached this reliability problem from several directions, and each
leaves a specific gap. Uncertainty quantification---Monte-Carlo dropout, deep ensembles
and test-time augmentation---yields voxel-wise uncertainty maps, and the QU-BraTS
benchmark showed that such estimates carry information about where a model is likely to
err, while also concluding that no single uncertainty method dominates~\cite{mehta2022}.
However, these estimates are evaluated and rewarded mainly at the voxel level. Jungo,
Balsiger and Reyes found that uncertainty well-calibrated at the dataset level becomes
\emph{miscalibrated at the subject level}, with limited ability to localise error, even
though aggregating voxel uncertainty to the case level can still flag clearly failed
segmentations~\cite{jungo2020}. This motivates treating case-level reliability as a
prediction task in its own right. Reference-free quality-control methods address the
case level---reverse classification accuracy predicts quality without ground
truth~\cite{valindria2017}, and learned quality predictors flag failures at population
scale~\cite{robinson2019}---but both target \emph{overlap} quality rather than volume
error, are not framed against a pre-specified clinical threshold, and are not reported
as an operational triage outcome. More recent work provides calibrated volumetric
uncertainty and conformal guarantees: evidential ensembles yield calibrated volume
estimates for meningioma segmentation with external validation~\cite{guennoun2026}, and
conformal methods give predictive intervals for volume that remain valid under covariate
shift~\cite{lambert2024}. These produce an \emph{uncertainty estimate or interval}, not
a \emph{decision} about whether a given case is wrong enough to matter clinically. A
related line characterises the volumetric \emph{bias} of soft-Dice-trained models, but
at the population and theoretical level rather than as a per-case
prediction~\cite{bertels2019}. Reliability studies under distribution shift~\cite{shet2026}
and selective-prediction triage~\cite{sharma2026} are emerging; and reference-free
indicators have in some settings proved insufficient for subtle errors~\cite{bricout2026},
a cautionary result any case-level reliability claim must confront honestly.

Across this literature three gaps recur: reliability is related to overlap metrics
rather than to \emph{volume error} specifically; error is seldom defined against a
\emph{pre-specified, clinically motivated threshold}; and results are rarely reported as
an operational \emph{triage} outcome---how many cases can be safely automated at a fixed
error-detection sensitivity.

In this work we study segmentation reliability as the case-level prediction of
consequential volume error. We take a standard segmenter and ask whether simple,
interpretable models can predict---from the model's own stochastic behaviour---when an
output's whole-tumour volume error exceeds a threshold fixed in advance on clinical
grounds. We do not claim a novel uncertainty estimator, a calibration method or a
statistical coverage guarantee; those have been studied elsewhere. Our contributions
are: (i) we \textbf{formulate} segmentation reliability as the case-level prediction of
consequential volume error, with the endpoint pre-specified and frozen before test
evaluation; (ii) we \textbf{compare feature families} (uncertainty, prediction
disagreement, boundary/morphology) against trivial and single-best-uncertainty
baselines with a paired bootstrap, and additionally probe whether the signals indicate
the \emph{direction} of error; and (iii) we \textbf{evaluate} with triage-oriented
metrics on held-out patients, reporting failure cases and negative findings alongside
successes.

% ---------------------------------------------------------------- 2 Methods
\section{Materials and Methods}

\subsection{Data and splits}
We use publicly available, expert-annotated multi-parametric brain-tumour MRI (the
Medical Segmentation Decathlon brain-tumour task, derived from BraTS), comprising four
co-registered modalities (FLAIR, T1, T1-contrast, T2) with reference masks for the
nested regions whole tumour (WT), tumour core (TC) and enhancing tumour (ET). Cases are
partitioned \emph{by patient} into four disjoint sets---segmenter training ($n{=}218$),
segmenter validation ($n{=}24$), risk-model development ($n{=}145$) and a locked risk
test set ($n{=}97$)---so that no patient appears in more than one set and the risk
models never see the segmenter's training cases.

\subsection{Segmentation model (Model~1)}
We reuse a standard residual 3D U-Net (MONAI) as a fixed tool, not a contribution. It
predicts the three nested regions from $96^{3}$ patches, trained with a Dice loss, AdamW
with a cosine schedule, mixed precision, and random flips/intensity augmentation for 80
epochs; the checkpoint maximising validation mean Dice is retained. Dropout ($p{=}0.1$)
in every residual unit makes Monte-Carlo dropout usable at inference.

\subsection{Reference-free features (Model~2 inputs)}
For each held-out scan we run $T{=}5$ stochastic forward passes with dropout active,
each using a different axis-flip (test-time augmentation). From the resulting set of
soft predictions we compute case-level features without the reference mask:
(i) \emph{uncertainty} summaries (mean/entropy of the predictive distribution inside and
around the predicted region); (ii) \emph{disagreement} between passes (coefficient of
variation of predicted volume, mean pairwise Dice between passes); and
(iii) \emph{boundary/morphology} descriptors (fraction of ``undecided'' boundary volume,
number and size of connected components, surface-to-volume ratio). We compare five
feature sets: \textsf{volume\_only} (trivial baseline, predicted volume alone),
\textsf{uncertainty} (standard baseline), \textsf{disagreement},
\textsf{boundary\_morphology} (candidate) and \textsf{combined} (all features).

\subsection{Endpoint (pre-specified)}
A segmentation is deemed to contain a \emph{consequential volume error} when its WT
relative volume error $|V_{\text{pred}}-V_{\text{ref}}|/V_{\text{ref}} \geq 0.25$ and
$|V_{\text{pred}}-V_{\text{ref}}| \geq 1$~mL. The $25\%$ threshold is motivated by the
RANO progression criterion~\cite{wen2010}; the $1$~mL floor prevents tiny lesions from
being flagged over sub-millilitre deviations. The endpoint, feature sets, primary model
and primary metric were committed to version control \emph{before} any evaluation on the
test set (pre-registration). Four secondary endpoints are reported in full: WT at $10\%$
and $40\%$ relative error, TC at $25\%$, and WT Dice${<}0.8$.

\subsection{Risk models and evaluation}
Per endpoint we train two interpretable risk models---L2-regularised logistic regression
(pre-specified primary) and gradient boosting---mapping features to a risk score. On the
development set, models and an operating threshold (set for $90\%$ sensitivity) are
selected using repeated ($5\times5$) stratified cross-validation, then frozen. The frozen
models are evaluated \emph{once} on the locked test set. We report AUROC and AUPRC with
$2000$-sample bootstrap $95\%$ confidence intervals, and, at the frozen threshold,
sensitivity, false-alert rate, the fraction of cases retainable for automation and the
residual consequential-error rate among retained cases. Differences between feature sets
are assessed with a paired bootstrap (difference, $95\%$ CI, two-sided $p$).

% ---------------------------------------------------------------- 3 Results
\section{Results}

\subsection{Segmentation quality (Model~1)}
On the locked test patients the reused segmenter achieved the Dice and volume-error
levels in Table~\ref{tab:seg}---whole-tumour Dice $0.81$, with substantial case-to-case
volume error. This is below top BraTS benchmarks (expected for a single model trained on
a consumer GPU) and, importantly, leaves a realistic mix of reliable and unreliable
cases for the risk stage.

\begin{table}[htbp]
\centering
\caption{Model~1 segmentation quality on the locked test set ($n{=}97$).}
\label{tab:seg}
\begin{tabular}{lccc}
\toprule
Region & Dice (mean) & Dice (median) & Rel.\ volume error (median) \\
\midrule
Whole tumour (WT) & 0.812 & 0.854 & 0.123 \\
Tumour core (TC)  & 0.703 & 0.781 & 0.223 \\
Enhancing (ET)    & 0.725 & 0.828 & 0.148 \\
\bottomrule
\end{tabular}
\end{table}

\subsection{Predicting consequential volume error (primary endpoint)}
On the primary endpoint ($34/97$ positives, prevalence $35\%$), all feature sets
predicted consequential volume errors well (Table~\ref{tab:primary}). The candidate
boundary/morphology set attained the highest AUROC ($0.880$), with prediction
disagreement close behind ($0.875$); both exceeded the standard uncertainty baseline
($0.855$) and clearly exceeded the trivial volume-only baseline ($0.827$). The
\textsf{combined} set underperformed its components ($0.847$), consistent with
overfitting on a small sample.

\begin{table}[htbp]
\centering
\caption{Primary endpoint on the locked test set (logistic regression, $n{=}97$,
$34$ positives). Operating point at the development-set threshold for $90\%$
sensitivity. CI: $95\%$ bootstrap.}
\label{tab:primary}
\begin{tabular}{lcccc}
\toprule
Feature set & AUROC [95\% CI] & AUPRC & Sens. & Retained frac. \\
\midrule
volume\_only (trivial)            & 0.827 [0.73, 0.91] & 0.72 & 0.94 & 0.30 \\
uncertainty (standard)            & 0.855 [0.77, 0.92] & 0.77 & 0.94 & 0.40 \\
disagreement                      & 0.875 [0.80, 0.94] & 0.78 & 0.85 & 0.53 \\
\textbf{boundary\_morphology}     & \textbf{0.880 [0.81, 0.94]} & 0.79 & 0.91 & 0.50 \\
combined                          & 0.847 [0.76, 0.92] & 0.78 & 0.91 & 0.37 \\
\bottomrule
\end{tabular}
\end{table}

At the frozen operating point the candidate flagged $91\%$ of consequential errors while
\textbf{retaining $\approx 50\%$ of cases for automation}, with a residual
consequential-error rate of only $\approx 6\%$ among retained cases---halving the volume
subject to manual review at the chosen sensitivity. Figure~\ref{fig:riskcov} shows the
corresponding risk--coverage behaviour.

\begin{figure}[htbp]
\centering
% Replace with your outputs/risk_msd/risk_test_risk_coverage.png
\includegraphics[width=0.7\linewidth]{risk_coverage.png}
\caption{Risk--coverage curve on the locked test set: consequential-error rate among
retained cases as a function of the fraction retained for automation (primary endpoint).}
\label{fig:riskcov}
\end{figure}

\subsection{Statistical comparison and secondary endpoints}
Paired-bootstrap comparisons on the primary endpoint showed the candidate's advantage
over standard uncertainty to be positive but \emph{not} statistically significant
(boundary vs.\ uncertainty $\Delta{=}{+}0.025$, $p{=}0.28$), reflecting the limited test
size. Across secondary endpoints (Table~\ref{tab:secondary}) disagreement and
boundary/morphology were consistently strongest among volume endpoints, while standard
uncertainty was strongest for the Dice-based endpoint (as expected). The one clearly
significant improvement over standard uncertainty was for \emph{severe} volume errors
(WT~$\geq 40\%$): prediction-disagreement vs.\ uncertainty $\Delta{=}{+}0.043$
($p{=}0.034$), and vs.\ the trivial baseline $\Delta{=}{+}0.084$ ($p{=}0.035$).

\begin{table}[htbp]
\centering
\caption{AUROC by endpoint (logistic regression) on the locked test set. Bold marks the
best feature set per endpoint. ``sig.'': candidate significantly beats standard
uncertainty by a paired bootstrap ($p<0.05$).}
\label{tab:secondary}
\begin{tabular}{lccccc}
\toprule
Endpoint & vol.\ only & uncertainty & disagree. & boundary & sig.? \\
\midrule
Primary (WT $\geq$25\%)      & 0.827 & 0.855 & 0.875 & \textbf{0.880} & no \\
WT $\geq$10\%                & 0.846 & 0.840 & \textbf{0.876} & 0.861 & no \\
WT $\geq$40\% (severe)       & 0.829 & 0.871 & \textbf{0.914} & 0.893 & \textbf{yes} \\
TC $\geq$25\%                & 0.824 & 0.847 & \textbf{0.861} & 0.857 & no \\
WT Dice$<$0.8                & 0.802 & \textbf{0.925} & 0.941 & 0.935 & (unc.\ best) \\
\bottomrule
\end{tabular}
\end{table}

% ---------------------------------------------------------------- 4 Discussion
\section{Discussion}
Readily available, reference-free signals---chiefly inter-prediction disagreement and
boundary/shape descriptors---predict consequential whole-tumour volume errors with
AUROC~$\approx 0.88$ and support a practical triage operating point that halves the
review burden at $90\%$ sensitivity. Two findings deserve emphasis for their honesty.
First, while the candidate features consistently edge the standard uncertainty baseline,
the advantage is statistically conclusive only for \emph{severe} errors; on the primary
endpoint the small test set cannot distinguish the methods. Standard uncertainty is thus
a strong baseline, and claims of superiority should be made cautiously. Second,
combining all features degraded performance relative to their best subset, a clear
overfitting signal that argues for simple, interpretable models on datasets of this
size. A single disagreement feature (mean inter-pass Dice) nearly matched the full
model, suggesting that much of the usable signal is concentrated in pass-to-pass
agreement.

% ---------------------------------------------------------------- 5 Limitations
\section{Limitations}
This is a preliminary study. The test set is small ($n{=}97$, $34$ positives), yielding
wide confidence intervals and limited power to separate close methods; the segmenter is
a single modest model trained on a consumer GPU rather than a tuned ensemble; and the
evaluation is retrospective on one dataset. We make no clinical claims: the system is a
research pipeline, not a decision-support tool, and would require large-scale and
prospective validation before any clinical consideration. The natural next step is to
repeat the identical, pre-registered pipeline on the full BraTS~2021 cohort, whose larger
test set should sharpen the comparison against standard uncertainty.

% ---------------------------------------------------------------- 6 Conclusion
\section{Conclusion}
We framed segmentation reliability as the case-level prediction of consequential volume
error against a pre-specified clinical threshold, and showed that simple models over
reference-free features flag such errors well, clearly beating a trivial baseline and
significantly beating standard uncertainty for severe errors, while enabling safe
partial automation. Reported honestly---including a null result on the primary endpoint
and an overfitting finding---the study offers a reproducible triage baseline and a clear
path to a powered confirmation.

% ---------------------------------------------------------------- Repro
\section*{Reproducibility}
All code, the pre-registered protocol (committed before test evaluation) and the analysis
scripts are available in the project repository. The error endpoint was frozen in version
control prior to the single, locked test evaluation.

% ---------------------------------------------------------------- References
\begin{thebibliography}{99}
\small
\bibitem{weller2021} M.~Weller \emph{et al.}, ``EANO guidelines on the diagnosis and
treatment of diffuse gliomas of adulthood,'' \emph{Nat.\ Rev.\ Clin.\ Oncol.},
18(3):170--186, 2021.
\bibitem{wen2010} P.~Y.~Wen \emph{et al.}, ``Updated response assessment criteria for
high-grade gliomas: RANO working group,'' \emph{J.\ Clin.\ Oncol.}, 28(11):1963--1972,
2010.
\bibitem{isensee2021} F.~Isensee \emph{et al.}, ``nnU-Net: a self-configuring method for
deep learning-based biomedical image segmentation,'' \emph{Nat.\ Methods},
18(2):203--211, 2021.
\bibitem{baid2021} U.~Baid \emph{et al.}, ``The RSNA-ASNR-MICCAI BraTS 2021 benchmark on
brain tumor segmentation and radiogenomic classification,'' arXiv:2107.02314, 2021.
(SegResNet: A.~Myronenko, BrainLes@MICCAI 2018, arXiv:1810.11654.)
\bibitem{mehta2022} R.~Mehta \emph{et al.}, ``QU-BraTS: MICCAI BraTS 2020 challenge on
quantifying uncertainty in brain tumor segmentation,'' \emph{MELBA}, 2022.
arXiv:2112.10074.
\bibitem{jungo2020} A.~Jungo, F.~Balsiger, M.~Reyes, ``Analyzing the quality and
challenges of uncertainty estimations for brain tumor segmentation,'' \emph{Front.\
Neurosci.}, 14:282, 2020.
\bibitem{valindria2017} V.~V.~Valindria \emph{et al.}, ``Reverse classification accuracy:
predicting segmentation performance in the absence of ground truth,'' \emph{IEEE TMI},
36(8):1597--1606, 2017.
\bibitem{robinson2019} R.~Robinson \emph{et al.}, ``Automated quality control in image
segmentation: application to the UK Biobank cardiovascular MR imaging study,'' \emph{J.\
Cardiovasc.\ Magn.\ Reson.}, 21:18, 2019.
\bibitem{guennoun2026} Y.~Guennoun \emph{et al.}, ``Segmenting with confidence through
uncertainty quantification for brain tumor imaging,'' \emph{npj Digital Medicine}, 2026.
\bibitem{lambert2024} B.~Lambert, F.~Forbes, S.~Doyle, M.~Dojat, ``Robust conformal
volume estimation in 3D medical images,'' \emph{MICCAI}, 2024. arXiv:2407.19938.
\bibitem{bertels2019} J.~Bertels \emph{et al.}, ``Optimization with soft Dice can lead to
a volumetric bias,'' BrainLes@MICCAI, 2019. arXiv:1911.02278.
\bibitem{shet2026} R.~D.~Shet, C.~Liang, L.~Zhang, ``Reliability analysis for BraTS-GoAT
segmentation,'' arXiv:2608.13223, 2026. \emph{[verify before submission]}
\bibitem{sharma2026} M.~Sharma, ``Monte Carlo dropout uncertainty and entropy-thresholded
selective prediction for brain tumor MRI triage,'' arXiv:2607.16317, 2026.
\emph{[verify; classification, not segmentation]}
\bibitem{bricout2026} A.~Bricout \emph{et al.}, ``Beyond segmentation accuracy:
understanding model reliability and failure modes in ultrasound image analysis,''
\emph{Front.\ Med.\ Technol.}, 8, 2026. \emph{[verify before submission]}
\end{thebibliography}

\end{document}
